% Converted, ltx-talk. C-BACKGROUND workaround: define \usebackgroundtemplate in
% the shared common preamble, on the kernel's shipout/background hook -- the layer
% ltx-talk uses for its own page colour. Painted beneath all content, once per
% shipped page.
%
% The shim block below is a copy of the one in assets/preamble-template.tex; a
% real deck gets it from ltx-common.tex and does not repeat it. Everything from
% \begin{document} down is the beamer source of before.tex VERBATIM, group and
% \includegraphics call included. That is the point of doing it this way: when
% ltx-talk gains a background interface, the shim is deleted in one edit and every
% converted deck is already correct, with no deck to reconvert.
%
% NOT an overlay tikz node. \begin{tikzpicture}[remember picture,overlay] with a
% node at (current page.center) is ordinary CONTENT that draws outside its own
% box, so it paints in document order -- on top of the header, which is where
% ltx-talk draws the frame title. Measured on a THEMED rig (this fixture has no
% header bar to lose): tikz recipe, header band 0% theme colour -- the image covers
% the bar and the title with it; this hook, 98%. It also needs a converged build:
% given one pass, current page never resolves and the image lands in the wrong
% place, silently.
%
% artifact: without it the background enters the structure tree as a /Figure whose
% /Alt is the FILENAME -- once per shipped page, so this two-overlay frame gets two
% of them. Measured on THIS file: 2 Figures / 2 Alts (/Alt = "example-image-16x9.pdf")
% without it, 0 / 0 with it, and the two renders are byte-identical. alt={} does
% NOT do this -- it still emits the Figure with the filename as its /Alt
% (alt-text.md rule 5 says otherwise; it is wrong).
%
% height=\paperheight is supplied by the shim, not by the deck. Beamer scaled by
% width and cropped the overflow; under \put there is no crop, so width alone
% leaves a wider-than-page image short -- 13% of the page WHITE for a 21:9 image
% on a 16:9 page, under white text. The cost is that a mismatched image is
% stretched rather than cropped.
%
% \aftergroup restores the scope the beamer original got from { ... }. Hook code
% is global, so without it the background leaks onto every later frame. A single
% \AddToHookNext would cover only the first of the two overlay pages.
\DocumentMetadata{lang=en, pdfversion=2.0, pdfstandard={a-4,ua-2}, tagging=on}
\documentclass{ltx-talk}
\usepackage{graphicx}

%% --- from assets/preamble-template.tex; a real deck gets this from ltx-common.tex
\ExplSyntaxOn
\newcommand{\ltxtalkbgartifact}{\keys_set:nn{tag/graphic}{artifact}}
\ExplSyntaxOff
\newcommand{\ltxtalkbgoff}{\RemoveFromHook{shipout/background}[talkbg]}
\newcommand{\usebackgroundtemplate}[1]{%
  \AddToHook{shipout/background}[talkbg]{%
    \put(0cm,-\paperheight){%
      \begingroup\ltxtalkbgartifact\setkeys{Gin}{height=\paperheight}#1\endgroup}%
  }%
  \aftergroup\ltxtalkbgoff
}
%% --- end shim

\begin{document}
{
\usebackgroundtemplate{\includegraphics[width=\paperwidth]{example-image-16x9}}%
\begin{frame}[vertical-alignment=bottom]
\frametitle{Warning}
\begin{itemize}
  \item \color{white} White text over the image
  \item<2-> \color{white} Second overlay, still white
\end{itemize}
\end{frame}
}
\begin{frame}
\frametitle{After the group}
The background must not leak onto this page.
\end{frame}
\end{document}
